\section{¿De qué bolsillo salen los ingresos de la IA? El fondo de ingresos electrónicos y el fondo laboral}

La sección anterior analizó cada segmento; esta responde una pregunta más de fondo: ¿de qué bolsillo salen los enormes ingresos de estas empresas de IA? Freda llama a los ingresos ya existentes de internet «ingresos electrónicos/impuesto electrónico»: publicidad en línea, comisiones de comercio electrónico y suscripciones, que en conjunto rondan los 400,000 millones de dólares. Este fondo no es pequeño, pero si una sola empresa, OpenAI, quisiera llegar a 200,000 millones de dólares de ingresos, obtenerlos únicamente de este fondo implicaría que Google, Meta, Amazon y otras cedieran una parte considerable de su participación.

\begin{figure}[H]
\centering
\includegraphics[width=0.94\textwidth]{electronic-revenue-vs-labor.png}
\caption{Ingresos electrónicos frente al fondo laboral: lo que ya existe en internet no es lo bastante grande; las cifras realmente grandes están en el trabajo y la productividad. Diagrama conceptual propio, elaborado a partir de la conversación entre 01:09:51 y 01:14:11.}
\end{figure}

\begin{knowledgebox}{Lectura de la figura: por qué disputarse solo la publicidad no basta}
A la izquierda, los ingresos electrónicos incluyen publicidad, comisiones de comercio electrónico y suscripciones; el programa los sitúa en el orden de 400,000 millones de dólares. A la derecha, el fondo laboral corresponde al costo de la mano de obra en Estados Unidos, del orden de 15 billones de dólares. Si la IA solo se convierte en «un Google pequeño», su alcance es limitado; solo si entra en la atención al cliente, el trabajo del conocimiento, AI for Science y el aumento de la productividad, el margen de ingresos crecerá de verdad.
\end{knowledgebox}

\subsection{¿Crecerán los ingresos publicitarios gracias a la IA?}

Freda no está de acuerdo con la idea de que «la IA hará crecer de forma notable el gasto publicitario total». Los ingresos publicitarios de Estados Unidos han crecido de manera relativamente estable en los últimos veinte años, no hay razón para que la proporción media de los ingresos que las empresas destinan a publicidad aumente de pronto y de forma marcada, y la penetración de la publicidad en línea ya es muy alta. Por ello, que la IA se dispute la publicidad se parece más a una redistribución de lo existente que a la creación, de la nada, de un gran fondo nuevo.

\begin{warningbox}{Hay que distinguir entre crecimiento publicitario y participación publicitaria}
La IA puede cambiar quién se queda con los ingresos publicitarios, pero no necesariamente amplía de forma notable el total del mercado publicitario. La competencia entre ChatGPT, Google, TikTok, Amazon y Meta es, sobre todo, una redistribución de participación y de puntos de acceso.
\end{warningbox}

\subsection{Fondo laboral, productividad e impuesto a la IA}

Si la capacidad de los modelos mejora, las cifras realmente grandes están en el costo de la mano de obra. El PIB de Estados Unidos ronda los 30 billones de dólares y el costo de la mano de obra, unos 15 billones; solo la atención al cliente es un mercado de cientos de miles de millones de dólares. Si la IA eleva un 10\% la productividad mundial, puede generar un enorme incremento del PIB. Freda lo ilustra con un ejemplo sencillo: 100 personas producen 100 bienes; con IA, 80 personas producen 110 bienes más baratos y mejores. El PIB real crece, y las utilidades de la empresa y los ingresos de los empleados que permanecen pueden aumentar, pero las 20 personas sustituidas pasarán por una transición dolorosa.

Por eso el programa también menciona la posibilidad de un impuesto a la IA o un impuesto a los despidos. No se trata de alarmismo, sino de un problema de distribución habitual en las revoluciones de la productividad: que la producción total crezca no significa que todos se beneficien a corto plazo.

\begin{importantbox}{Las dos etapas de los ingresos de la IA}
Al principio, cuando la capacidad de los modelos todavía no ha madurado del todo, a la IA le resulta más fácil disputarse los ingresos electrónicos: publicidad, suscripciones, comisiones de comercio electrónico y presupuestos de software. Más adelante, si los Agent y AI for Science maduran, la IA podrá entrar de verdad en el fondo laboral y en el fondo de productividad adicional.
\end{importantbox}

\subsection{Resumen de la sección}

Los ingresos de la IA provienen de dos niveles. Los ingresos ya existentes de internet pueden sostener la comercialización inicial, pero no bastan para explicar inversiones del orden del billón de dólares; las cifras realmente grandes están en el costo de la mano de obra, el aumento de la productividad y la creación de nuevo valor. Para juzgar si hay una burbuja, hay que ver si la IA logra pasar del primer nivel al segundo.
